% Inclusion fragment. Requires amsmath, amssymb, amsthm, and hyperref.
% Uses theorem, lemma, corollary, and remark environments supplied by the parent.
% Bibliography key required: pinsker-almost-unary.
% All source-size and interpretation claims must retain the stated scope.

\section{A stronger obstruction to removing the initializer}
\label{sec:finite-dimensional-obstruction}

The supplied transient-seed draft proves that every term of a particular
initializer-tied recursion language is bounded by a function of at most
one input coordinate. It consequently excludes ordinary addition and
pairing \emph{as numerical terms}. Its final scope paragraph expressly
leaves open changed domains, finite tuples, and quotient interpretations.
The present section strengthens that conclusion: arbitrary finite tuples
and arbitrary quotient equality still do not help when the interpreted
operation has one fixed tuple of term realizers. The gap left by this
particular argument is genuinely relational interpretation, rather than
merely a different finite-coordinate encoding.
Section~\ref{sec:relational-obstruction} later excludes relational
bi-interpretations for a specified ordinary tied-seed theory by an
additional proof-theoretic argument.

The inherited result is in
\file{ipra_transient_seed_2026-11-10.zip},
member \file{theory/obstruction.tex}, section
``A bounded-coordinate obstruction for a tied-seed recursor,'' especially
its ``Exact logical scope'' subsection. The counter-congruence draft
reproduces the same invariant in
\texttt{theory/obstructions.tex}, section
``A more promising seedless operator still has only one unbounded input.''
Both archives are direct members of the pinned ProveIt snapshot
\texttt{3fb35a5f2c0cb506180bede7ee3b9b5f93d4bc9d}.
The one-coordinate invariant is inherited; the finite-power quotient
argument below is the strengthening relative to those inspected drafts.
No literature-wide priority claim is made. The broader class of functions
bounded in one coordinate is conventionally called the class of
\emph{almost unary functions}; see Pinsker~\cite{pinsker-almost-unary},
Section~1.1.

\subsection{The exact numerical grammar being obstructed}

Terms are generated by
\[
             t::=x\mid 0\mid B(n;u,v.h),
\]
where only the update body $h$ binds the distinct variables $u,v$.
Updates are fixed finite earlier terms, with arbitrary passive numerical
parameters. Put
\[
             Sa:=B(0;u,v.a)
\]
when $u,v$ do not occur in $a$. The proposed computation schemas are
\begin{align}
 B(0;u,v.h)&=S\bigl(h[0/u,0/v]\bigr),\label{obs:B0}\\
 B(Sn;u,v.h)&=h[n/u,B(n;u,v.h)/v].\label{obs:BS}
\end{align}
Thus the initial state is tied to the same update used in the recursive
step; there is no independent initializer field. For each fixed body,
the standard natural-number interpretation is
\[
 B_h(0,\vec p)=1+h(0,0,\vec p),\qquad
 B_h(m+1,\vec p)=h(m,B_h(m,\vec p),\vec p).
 \tag{B$_{\mathbb N}$}\label{obs:B-standard}
\]
Structural recursion on descriptions, followed by numerical primitive
recursion, defines this interpretation. In particular the abbreviation
$S$ denotes ordinary successor. No uniform evaluator for all descriptions
is postulated.

\begin{lemma}[Inherited coordinate bound, with proof]
\label{obs:coordinate-bound}
For every fixed $B$-term $t(x_1,\ldots,x_r)$, either its standard values
are bounded by a constant, or there are an index $j$ and a nondecreasing
unary primitive-recursive function $F$ such that
\[
        t(x_1,\ldots,x_r)\le F(x_j)
        \quad\hbox{for all }(x_1,\ldots,x_r)\in\mathbb N^r.
        \tag{CB}\label{obs:CB}
\]
The coordinate and bound can be obtained from the finite syntax. The
bounding function is not asserted to be a term of the $B$ language.
\end{lemma}

\begin{proof}
Use external structural induction on terms, treating the two bound
variables of an update as independent input coordinates when applying
the induction hypothesis to its body. Variables and zero satisfy the
claim. Consider $B(n;u,v.h)$, renaming $u,v$ away from the free variables
of $n$ and the passive parameters.

If $h$ is globally bounded by $K$, the initial state is at most $K+1$
and every later state is at most $K$. If its controlling coordinate is a
passive parameter $p_j$, with $h(u,v,\vec p)\le H(p_j)$, every state is
at most $1+H(p_j)$. These bounds are independent of the counter value.

If its controlling coordinate is $u$, then $h(u,v,\vec p)\le H(u)$
for a nondecreasing $H$. Equation~\eqref{obs:B-standard} bounds the
state at counter $m$ by $1+H(m)$: this includes the base state, and the
positive state at $m+1$ is at most $H(m)$.

Finally, if its controlling coordinate is $v$, set
\[
 K(0)=1+H(0),\qquad
 K(m+1)=\max\{K(m),H(K(m))\}.
\]
This is a nondecreasing primitive-recursive function. External induction
on the standard counter proves $B_h(m,\vec p)\le K(m)$, uniformly in
all passive parameters: the base is immediate, and the next state is at
most $H(B_h(m,\vec p))\le H(K(m))\le K(m+1)$.

In the last two cases write the nondecreasing counter bound as $K$.
Apply the structural induction hypothesis to the actual counter term
$n$. A constant bound $n\le c$ gives the constant bound $K(c)$. A bound
$n(\vec x)\le N(x_j)$ gives $K(N(x_j))$. Composition preserves
monotonicity and primitive recursiveness. All possible controlling
coordinates of $h$ have now been covered.
\end{proof}

This is a theorem about standard term operations, not an object-language
induction principle. Adding axioms that are sound in this standard
algebra cannot change those term operations or invalidate the bound.

\subsection{The finite-power quotient theorem}

The next theorem uses only the conclusion of
Lemma~\ref{obs:coordinate-bound}; it applies to any algebra on
$\mathbb N$ whose term operations have that property. Its sets and
equivalence relations need not be definable.

\begin{theorem}[No infinite quotient with a term-defined two-sided unit]
\label{obs:quotient-unit}
Let $d\ge1$, let $D\subseteq\mathbb N^d$, let $\sim$ be any equivalence
relation on $D$, and let $e\in D$. Suppose that a map
\[
 \mu=(\mu_1,\ldots,\mu_d):D\times D\longrightarrow D
\]
is represented by one fixed tuple of terms in its $2d$ input coordinates,
with any fixed numerical parameters, and that
\[
             \mu(e,x)\sim x\sim\mu(x,e)
             \qquad(x\in D).
             \tag{Unit}\label{obs:unit}
\]
Then $D/{\sim}$ is finite.

More precisely, a composition of at most
\[
                  m=\left\lceil\log_2(d+1)\right\rceil
\]
maps of the forms $x\mapsto\mu(e,x)$ and $x\mapsto\mu(x,e)$ has finite
image and sends every element to its own $\sim$-class.
\end{theorem}

\begin{proof}
Define
\[
                    T_0(x)=\mu(e,x),\qquad T_1(x)=\mu(x,e).
\]
By closure and~\eqref{obs:unit}, both maps send $D$ into $D$ and preserve
each $\sim$-class. The same holds for every finite composition. Only
transitivity of $\sim$ is used here.

For each output coordinate $i$, choose a bound supplied by
Lemma~\ref{obs:coordinate-bound} for $\mu_i$. A bound controlled by one
of the fixed numerical parameters is a constant after those parameters
are fixed, so it can be included in the constant case. Mark a coordinate
with a constant bound as \emph{terminating}.

Every remaining coordinate has a controlling variable from one of the
two input tuples. If
\[
                 \mu_i(x,y)\le F_i(x_j),
\]
record an edge $i\longrightarrow j$ labelled $1$. In that case
\[
                 (T_0(x))_i\le F_i(e_j),\qquad
                 (T_1(x))_i\le F_i(x_j).
                 \tag{First input}\label{obs:first-input}
\]
If instead
\[
                 \mu_i(x,y)\le F_i(y_j),
\]
record an edge $i\longrightarrow j$ labelled $0$. Then
\[
                 (T_1(x))_i\le F_i(e_j),\qquad
                 (T_0(x))_i\le F_i(x_j).
                 \tag{Second input}\label{obs:second-input}
\]
Thus at every nonterminating coordinate there is exactly one label that
can leave dependence on the variable input in the selected bound, and
exactly one next coordinate to which that dependence passes. Choosing
the other label gives a constant bound immediately.

Starting at a fixed coordinate $i$, this deterministic labelled graph
admits at most one binary word of length $m$ along which a dependence
bound survives all $m$ steps. A path meeting a terminating coordinate
does not admit such a word. Consequently at most $d$ of the $2^m$ binary
words are admitted from any of the $d$ starting coordinates. Since
$2^m>d$, choose a word $w=b_1\cdots b_m$ admitted from none of them, and
put
\[
                T_w=T_{b_1}\circ T_{b_2}\circ\cdots\circ T_{b_m}.
                \tag{Composition order}\label{obs:word-order}
\]
For clarity, a coordinate bound is traced from the \emph{outermost}
operation to the innermost, so the word in~\eqref{obs:word-order} is read
from $b_1$ to $b_m$. For instance, if the selected bound for the outer
coordinate permits $b_1$ and leads from $i$ to $j$, then
\[
 (T_w(x))_i\le F_i\bigl((T_{b_2}\circ\cdots\circ T_{b_m}(x))_j\bigr).
\]
The choice of $w$ ensures that this chain reaches a terminating bound or
a mismatched edge label no later than the last letter. The bound at
that stage is a fixed constant. Composing outward with the preceding
nondecreasing unary functions yields a fixed constant $C_i$ independent
of $x$. This verifies the bound in the specified composition order; no
commutation or rearrangement of the maps $T_0,T_1$ is assumed.

Therefore
\[
           T_w(D)\subseteq
           D\cap\prod_{i=1}^{d}\{0,1,\ldots,C_i\},
\]
which is finite. Yet $T_w(x)\sim x$ for every $x\in D$. Every
$\sim$-class has a representative in this finite image, so the quotient
is finite.
\end{proof}

\begin{remark}[Very little algebra is needed]
The theorem does not use associativity or commutativity of $\mu$.
It does not even require $\mu$ to respect $\sim$ in both inputs; the
two displayed unit identities and closure suffice. In applications
where $\mu$ interprets a monoid operation on a quotient, compatibility
is available but unnecessary to this argument. If $d=0$, the ambient
power is a singleton and finiteness is immediate.
\end{remark}

\subsection{Consequence for interpretations of arithmetic}

\begin{corollary}[Finite dimension and quotient equality do not repair a term compiler]
\label{obs:no-ipra-term-interpretation}
Let $T$ be a theory satisfied by the standard $B$-algebra described in
\eqref{obs:B-standard}. There is no interpretation of full first-order
$\mathsf{iPRA}$ in $T$ having all the following properties:
\begin{enumerate}
\item the interpreted numerical domain is a nonempty definable subset
      $D$ of one finite power of the source domain;
\item interpreted equality is any definable equivalence relation on $D$;
\item interpreted addition has one fixed tuple of $B$-term realizers
      mapping $D\times D$ into $D$.
\end{enumerate}
The claim allows fixed parameters and any interpretation of successor
consistent with the target axioms. In particular, there is no
bi-interpretation meeting these restrictions.
\end{corollary}

\begin{proof}
Apply the assumed interpretation in the standard $B$-model of $T$.
Choose a representative $e$ of interpreted zero. The translated
addition unit laws give~\eqref{obs:unit}. By
Theorem~\ref{obs:quotient-unit}, the interpreted numerical quotient is
finite. But the translated target axioms give an injective successor
map on that quotient whose image omits zero. An injective self-map of a
finite nonempty set is surjective, a contradiction.
\end{proof}

The qualifier concerning the standard $B$-model is essential. A theory
whose additional axioms deliberately exclude that model is not ruled
out solely by an analysis of its standard term functions. The natural
candidate based on the displayed computations and rules sound over
$\mathbb N$ does have this model.

Likewise, the restriction on the interpreted operation is essential.
An ordinary first-order interpretation may define an operation by a
formula with a unique output modulo the interpreted equality, rather
than by a fixed tuple of source terms. Such graphs remain outside this
theorem. Even a finite definition by cases can choose different
controlling input coordinates in different cases. For example, the
operation $\max(x,y)$ on $\mathbb N$ has identity zero and is defined by
cases selecting one of its arguments, whereas no single coordinate
bounds it. This example explains the scope limitation; it is not a
claim that order or maximum is definable in the $B$-algebra.

The logarithmic bound in Theorem~\ref{obs:quotient-unit} counts layers of
unit multiplication in a metatheoretic construction. It is not a bound
on expanded numerical term size, evaluation time, or proof-checker
source size. The theorem strengthens a negative result for one proposed
smaller arithmetic basis. It does not provide the smaller, complete
bi-interpretable axiomatization requested in the task.
